\label{sec:loading}

Loading a file is part of program execution in ABCL/R2. Its expressions
run in an environment, using a group and evaluator; the same translated
file can therefore behave differently when loaded by different callers.
First separate preparing a file from executing it:

\par\medskip\noindent
\begin{minipage}{\linewidth}
\begin{tabular}{L{0.32\textwidth}L{0.58\textwidth}}
\textbf{Expression} & \textbf{Purpose}\\\hline
\code{(inc "name")} & Translate if needed, then load and evaluate in the
caller's environment, group, and evaluator context.\\
\code{(comp "name")} & Translate and compile if needed; does not run the
ordinary ABCL application expressions.\\
\code{(inc-lib "name")} & Load and evaluate in the global environment,
top-level group, and current top-level evaluator.\\
\code{(require 'name)} & Search for and load an unregistered module through
\code{inc-lib}.\\
\code{(provide 'name)} & Record a module name in \code{*modules*}; does not
load a file or validate its definitions.\\
\code{(pwd)} / \code{(cd "dir")} & Historical directory commands;
not implemented for SBCL in this port.\\
\code{:q} / \code{:reset} & Quit the interpreter or reset the system.\\
\end{tabular}
\end{minipage}\par\medskip

Use a file's \emph{base pathname}, without its type:
\code{(inc "demo" :use-source t)} loads \filepath{demo.abclr2}.
Translation writes a neighboring \filepath{demo.lsp}; compilation writes
a host-specific compiled file, normally \filepath{demo.fasl} on SBCL.
Do not pass the \code{.abclr2} type as part of the base name: the loader
derives all three paths from that argument.
The directory must permit writing the generated files. Use scratch
copies when exploring the original examples, as the checks do.

Relative paths are resolved against the process's current directory,
not automatically against the containing source file's directory.
Start the interpreter in the required directory or use absolute paths.
The host-specific implementations in \filepath{src/topfunc.lsp} have
no SBCL branch for changing that directory. Calling \code{inc-lib} does
not make its relative path relative to \filepath{lib/}.

\subsection{Options and File Selection}

The user's guide specifies the public options~\cite{guide}; the active
implementations are in \filepath{src/inc.lsp} and \filepath{src/user.lsp}.
The defaults are:

\noindent\begin{minipage}{\linewidth}
\begin{longtable}{L{0.20\textwidth}L{0.15\textwidth}L{0.55\textwidth}}
\textbf{Option} & \textbf{Default} & \textbf{Accepted by and effect}\\\hline
\code{:use-source} & \code{nil} & \code{inc}, \code{inc-lib}, \code{require},
\code{comp}. Non-nil forces translation from ABCL source. On a load it
also bypasses the compiled-file branch; on \code{comp} it forces compilation.\\
\code{:verbose} & \code{t} & All four commands. Controls translation and
the passed-through load/compile verbosity, not every warning or diagnostic.\\
\code{:compile} & \code{nil} & \code{inc}, \code{inc-lib}, \code{require}.
Non-nil runs \code{comp} before loading. \code{comp} itself has no
\code{:compile} option.\\
\end{longtable}
\end{minipage}

\code{inc} and \code{inc-lib} can prompt for a filename when it is
omitted. \code{comp} requires a name, and \code{require} requires a
module name. In scripts and reproducible sessions, always supply it
explicitly. \code{provide} takes only its module-name argument.

\paragraph{Follow the three-stage choice.}
For an explicit \code{inc} or \code{inc-lib}, read the loader as follows:

\begin{enumerate}
  \item If \code{:compile} is true, first perform the compilation command
        with the selected \code{:use-source} and \code{:verbose} values.
  \item Require the \code{.abclr2} source to exist. Translate it if forced,
        if there is no \code{.lsp}, or if the source is newer than the
        translation. Even a cached compiled file does not remove this
        source-file requirement.
  \item Load the compiled file only when \code{:use-source} is false,
        that file exists, and it is \emph{strictly newer} than the Lisp
        translation. Otherwise load the translation. Execute its
        expressions in the selected loading context.
\end{enumerate}

\code{comp} uses the same translation test, then compiles if forced,
if the compiled file is absent, or if the translation is newer than
the compiled file. Otherwise it reports that the compiled file is up
to date. It does not execute an ordinary ABCL object definition or
application send merely by compiling it. Translation still processes
macros; do not equate compilation with an absence of all host-side effects.

Timestamp comparisons are not content comparisons. Equal timestamps
do not satisfy ``strictly newer'' in the load test, so loading may select
the Lisp file even after compilation. The combination
\code{:compile t :use-source t} compiles first but still follows the
forced-source load branch. Compilation is not a promise that this
particular load executes the FASL. This is usually a performance/file
selection distinction, not a different intended program result.

For day-to-day editing, prefer the explicit source path:

\begin{codeblock}
\begin{lstlisting}
(inc "demo" :use-source t)          ; retranslate, load, and run
(comp "demo" :use-source t)         ; prepare compiled output; do not run
(inc "demo" :compile t)             ; update as needed, then load and run
\end{lstlisting}
\end{codeblock}

These are three independent commands, not a required sequence. Do not
repeat a loading command solely because a translation was cached:
\code{inc} and \code{inc-lib} execute the file's expressions on each call,
including assignments, object creation, and application sends.
\code{:verbose nil} can reduce output while keeping those effects.

\subsection{Execution Context and Module Loading}

\paragraph{An included fragment sees its caller's environment.}
The supplemental \filepath{tutorial/examples/15-loaded-fragment.abclr2}
contains just these two expressions:

\wholelisting{tutorial/examples/15-loaded-fragment.abclr2}

It deliberately depends on \code{load-token}; it is a fragment to load
in a prepared environment, not a self-contained program. Try:

\begin{codeblock}
\begin{lstlisting}
[load-token := :global]
[loaded-value := :unchanged]
[loaded-object := nil]
(let ((load-token :local) (loaded-value nil) (loaded-object nil))
  (inc "tutorial/examples/15-loaded-fragment" :use-source t)
  (list loaded-value [loaded-object <== :value]))
; => (:LOCAL :LOCAL)
loaded-value                       ; => :UNCHANGED
\end{lstlisting}
\end{codeblock}

The fragment updates the local \code{loaded-value}, binds its new object
locally, and that object's script retains the local \code{load-token}.
The same fragment loaded by \code{inc-lib} instead uses the global
environment:

\begin{codeblock}
\begin{lstlisting}
(let ((load-token :local) (loaded-value nil) (loaded-object nil))
  (inc-lib "tutorial/examples/15-loaded-fragment" :use-source t)
  (list loaded-value loaded-object)) ; => (NIL NIL)
loaded-value                        ; => :GLOBAL
[loaded-object <== :value]           ; => :GLOBAL
\end{lstlisting}
\end{codeblock}

Group context matters too. Ordinary objects created by \code{inc} inside
a client's script inherit that caller's group; \code{inc-lib} evaluates
the file in the top-level group using the current \code{*eval*}.
An explicit \code{group} or \code{meta-gen} in the loaded definitions
can still choose different machinery. Loading a library does not itself
replace existing objects' evaluator references or compose its policy
with every custom group. This is why the original examples' setup and
loading order must be read as part of their protocols.

\paragraph{Require a module name, not an arbitrary file path.}
\label{sec:library-path}
\code{*library-directories*} contains the directories searched by the
ABCL/R2 \code{require} macro. It selects the first directory containing
the named \code{.abclr2} source, then uses \code{inc-lib} to load it.
The default value is \code{(list "")}: only the current working directory
is searched. The distribution's \filepath{lib/} directory is not added
automatically. Before running examples that require its modules, add it
at the ABCL/R2 prompt, starting the interpreter at the repository root:

\begin{codeblock}
\begin{lstlisting}
[*library-directories* :=
  (cons #P"lib/" *library-directories*)]
(require 'future)
\end{lstlisting}
\end{codeblock}

Here \code{\#P"lib/"} is Common Lisp's pathname literal for a directory.
Keep the trailing slash: the search uses the directory component of
each entry. Relative entries are resolved against the working directory,
not the executable's location or the directory of the file being loaded.
Use an absolute directory pathname when starting elsewhere.
Explicit \code{inc} and \code{inc-lib} calls do not search this list;
\code{(inc-lib "lib/future")} instead names the file directly.
The path is session configuration; after \code{:reset}, set it again.

Symbols are converted to lowercase names. String arguments retain their
case during search; prefer symbols or explicitly lowercase strings.
\code{provide} normalizes both symbols and strings to lowercase.

For your own library directory, use a directory pathname, save the
previous search path, and restore it after the experiment:

\begin{codeblock}
\begin{lstlisting}
[saved-library-directories := *library-directories*]
[*library-directories* :=
  (cons #P"/absolute/path/to/my-libraries/" *library-directories*)]
; (require 'my-module) ; loads my-module.abclr2 from the search path
[*library-directories* := saved-library-directories]
\end{lstlisting}
\end{codeblock}

The absolute example directory is a placeholder to replace, not a
required installation path.
Putting it first makes its source win over later directories. A typical
module includes \code{(provide 'my-module)} to register its name in
\code{*modules*}. \code{require} does not register a module automatically:
without \code{provide}, a second require can execute the file again.

Once a name is registered, \code{require} skips loading entirely, even
with \code{:use-source t} or \code{:compile t}. To intentionally reload
a known file while editing, use its explicit base pathname with
\code{inc-lib}. Be aware that this reexecutes definitions and effects;
it is not an atomic replacement of every reference to earlier objects.
If \code{provide} ran before a later failure in a module, its recorded
name alone is not evidence that all initialization succeeded.

A missing module also returns \code{nil} without loading in this
implementation. That return value is not a success certificate or a
reliable distinction between ``already registered'' and ``not found.''
Inspect the transcript, expected definitions, and search path. Direct
\code{inc} and \code{inc-lib} instead signal an error for missing source.

Included libraries are \code{future}, \code{multicast}, \code{express},
\code{heavy-meta}, \code{copy}, \code{watch}, and \code{report}.
Their walkthroughs explain setup and use. Historical alternatives with
shared feature or generator names require separate fresh sessions, not
an assumption that \code{require} can distinguish filenames by itself.

\paragraph{Evidence and check your understanding.}
Run \code{make -C tutorial loading-check}. Thirty-five probes check
option defaults, compile-versus-load effects, independently marked Lisp
and FASL cache selection, forced source, missing-source rejection,
module registration/reload/search, and local versus global environment,
group, and evaluator behavior. Tests modify only generated cache files
in \filepath{build/loading-tests/}; original sources remain unchanged.
The small fragment above is copied there before testing.

Why might changing a module and calling \code{require} show no change?
\emph{Answer:} The registered name bypasses the load, including its
source option. Why can \code{comp} create output without creating the
application object? \emph{Answer:} Preparation does not execute its
ordinary ABCL expressions. Why does a file see a local value with
\code{inc} but the global value with \code{inc-lib}?
\emph{Answer:} Loading selects an execution environment, not just a file.
